\documentclass[12pt,a4paper]{article}

% Encoding
\usepackage[utf8]{inputenc}
\usepackage[T1]{fontenc}

% Page layout
\usepackage{geometry}
\geometry{a4paper,left=2.54cm,right=2.54cm,top=2.54cm,bottom=2.54cm}

% Font - Arial (LSE requirement)
\usepackage{helvet}
\renewcommand{\familydefault}{\sfdefault}

% Spacing
\usepackage{setspace}

% Bibliography
\usepackage[
backend=biber,
style=apa,
url = false,
doi = false,
isbn=false,
]{biblatex}
\addbibresource{bibliography.bib}
\usepackage[hidelinks]{hyperref}

% Headers
\usepackage{fancyhdr}
\setlength{\headheight}{14.5pt}
\pagestyle{fancy}
\fancyhf{}
\rfoot{\thepage}
\renewcommand{\headrulewidth}{0pt}

% Fix hbox warning
\sloppy

% Misc
\usepackage{graphicx}
\usepackage{booktabs}
\usepackage{appendix}

\begin{document}

\begin{doublespace}

% Cover page
\thispagestyle{empty}
\pagenumbering{gobble}
\begin{center}
    \vspace*{2cm}
    {\large GV499 Dissertation\\
    MSc Public Policy and Administration \\
    Department of Government\\
    London School of Economics}\\
    \vspace{2cm}
    {\Large\textbf{Assumptions Matter: Central Bank Independence, Epistemology, and the Foreclosure of Democratic Contestation}}\\
    \vspace{2cm}
    {\large \textit{Dissertation Outline}}\\
    \vspace{0.5cm}
    \today
\end{center}
\newpage

\pagenumbering{arabic}
\setcounter{page}{1}

% -------------------------------------------------------

\section*{To be submitted as brief description of agreed topic on Topic Approval Form}

\section*{Research Question}
Why are the goals of monetary policy removed from democratic contestation, and what role does the economics literature on central bank independence play in that removal?

\section*{Thesis}
Central bank independence, as constructed in the economics literature, folds the \textit{goals} of monetary policy into the \textit{definition} of independence. This forecloses democratic contestation of those goals: to revise the goal is, by definition, to reduce independence. The assumptions behind this construction originate in the theoretical economics literature (the ``conservative central banker''), are operationalized in the empirical measurement of de jure central bank independence, and are transmitted through personnel and discourse into the consistent blockage of legislative revision of the Fed's dual mandate. Because monetary policy has real distributional consequences, this foreclosure is a substantive political choice disguised as a technical one --- and the goal it locks in happens to be the creditor-friendly conservative-banker default.

The broader point: the economics literature that justifies central bank independence removes the goals of monetary policy from democratic contest by defining independence as adherence to a particular goal. These papers have had an important epistemological influence on policymaking.

\section*{Methodology}

The argument will be constructed with a theory-testing single case study of the Federal Reserve, using the process tracing methodology described by Beach and Pedersen (2019). I build my causal mechanism through three linked points: the epistemological prior from economics literature folds the goal of monetary policy into the definition of independence; those assumptions travel into policy and institutional decisions; that transmission forecloses legislative revision. Each step is downstream of the prior one. To build this argument, I will reference the assumptions built into economics literature, along with public Federal Reserve documents, like speeches and policy statements.


\newpage

\textbf{Notes for Cheryl Meeting, 5-11}

Make sure to be super precise about the CBI that I'm referring to: goal, target, and operational

Don't overstep; stick to the emperics and improving it

Explain why process tracing is defensible

Reach out to Adam Posen

Send Cheryl an email with a few of these questions- regarding research

``There are academics who like to jump on a bandwagon of critiquing central banks"

Citations; neutrality of money, breaking down independence forms,

Give it some thought on where exactly I am re-examining CBI

Write as an academic paper not as a blog

\newpage

\tableofcontents
\newpage

% -------------------------------------------------------
\section{Introduction (\textasciitilde750 words)}

\begin{itemize}
    \item \textbf{HOOK: JPMorgan/Powell paradox} --- JPMorgan defends Powell against a government it otherwise influences.

    \item \textbf{The puzzle} --- CBI is treated as apolitical and technical, yet it determines who bears the costs of monetary policy. Congress formally has control over the goals of the Federal Reserve. Why is the question of those goals absent from democratic debate?

    \item \textbf{RQ + thesis preview}

    \item \textbf{Epistemic foreclosure as central contribution} --- The economics literature that defends the merits of CBI folds the goal of monetary policy into the definition of independence, removing the goal from democratic contest by never posing it as a question.

    \item \textbf{``Ingrained'' framing} --- The normative default was built into the epistemological framing of CBI from the start; it did not have to be lobbied for.
\end{itemize}

% -------------------------------------------------------
\section{Background and Theory (\textasciitilde1,900 words)}

\subsection{The Economics Literature on CBI (\textasciitilde700 words)}

\subsubsection{Central banking and CBI defined}

\begin{itemize}
    \item \textbf{Goal vs.\ instrument independence (Debelle and Fischer 1994)} --- Congress sets the goals; the Fed has operational independence to achieve those goals. Powell himself says this. The official position is that independence is instrument-only. This distinction matters for the analysis of CBI indices, where the price stability component is nonetheless incorporated into CBI scores.

    \item \textbf{Fed's dual mandate (Section 2A, Federal Reserve Act)} --- Simple definition. Worth noting this is actually a ``tripartite'' mandate.
\end{itemize}

\subsubsection{Theoretical foundations}

\begin{itemize}
    \item \textbf{Kydland and Prescott (1977)} --- Time-inconsistency problem applied to the inflation-unemployment tradeoff. The general principle for CBI.

    \item \textbf{Barro and Gordon (1983)} --- Applies the logic directly to monetary authorities. Monetary policy is neutral in the long run, so it can only reap benefits through short-run surprise. Under discretion, this produces an inflationary bias.

    \item \textbf{Rogoff (1985)} --- The conservative central banker solution: appoint a monetary authority who prefers less inflation than society does.

    \item \textbf{Posen (1995)} --- The conservative central banker's preferences resemble those of the financial sector; financial-sector support is associated with levels of CBI. \textit{Scope: this is a claim about whose preferences the locked-in default resembles, not a causal claim about who built it.}
\end{itemize}

\subsubsection{Empirical development}

\begin{itemize}
    \item \textbf{GMT (1991); CWN (1992); Alesina and Summers (1993)} --- Observable institutional characteristics compiled into indices; statistical models measuring how CBI influences output and inflation. The free lunch claim: lower inflation, no output cost. Going to address price stability conservatism in \ref{sec:epistemology}.

    \item \textbf{Romelli (2022); Adrian et al.\ (2024)} --- Same general approach, updated and extended. Romelli's objectives dimension --- which encodes price stability conservatism --- has more than doubled across his sample period.

    \item \textbf{Shared features across all five} --- De jure measurement, inherited time-inconsistency prior, price stability as the benchmark, free lunch claim that drops Rogoff's acknowledged trade-off. They don't argue against employment as an objective --- they assume it away by construction.

    \begin{itemize}
        \item There isn't a great solution to empirical measurement of CBI beyond de jure indices
    \end{itemize}

    \item \textbf{Isabel Schnabel on monetary dominance} --- Notes that the dual mandate means little in practice. Monetary dominance is (a) reflected in the current conduct of monetary policy and (b) puts the price stability mandate first. Big political science implications

\end{itemize}

\subsubsection{The technocratic legitimation claim}

\begin{itemize}
    \item \textbf{Fischer (1995)} --- Instrument independence justified by technical complexity. Monetary policy is too complex for democratic deliberation; it requires expert insulation.

    \item \textbf{Does double work} --- Justifies CBI on meritocratic grounds AND pre-emptively delegitimises democratic oversight of the goal.

    \item \textbf{Flag} --- Bernanke in Section 4.3 confirms the informational asymmetry this claim relies on. Posen/Cochrane (2025) essentially restate this claim as well.
\end{itemize}

% -------------------------------------------------------
\subsection{The Political Science Literature on CBI (\textasciitilde400 words)}

\textit{Establish two core facts: (1) There is a non-engaged-with political science CBI literature, particularly regarding the democratic implications of monetary policy. (2) Central banks have a unique relationship with finance.}

\subsubsection{Literature on CBI and non-neutrality of central banks/bankers}

\begin{itemize}
    \item \textbf{Adolph (2018)} --- Shadow principals: career incentives shape central banker preferences before and after office. Greater legal independence makes central bankers more powerful agents and more valuable to shadow principals. Critiques the ``myth of neutrality.''

    \item \textbf{Wansleben (2020)} --- The Fed as bureaucracy and institution building

    \begin{itemize}
        \item \textbf{Mabbett and Schelkle (2019)} --- Broader discussion of the Fed's ``loneliness.''
    \end{itemize}

    \item \textbf{Jacobs and King (2018)} --- The Fed has historically justified growing its powers through technocratic merit.

    \item \textbf{Downey (2025)} --- Double delegation of monetary powers.

    \item \textbf{Zingales (2013)} --- Economists face the same capture mechanisms as regulators: career incentives, funding dependencies, professional socialisation. Pro-market findings circulate more widely and define what is acceptable to conclude.

    \item \textbf{Federal Advisory Council finding} --- Only the four largest banks on the FAC demonstrably influence the Federal Funds rate (these banks matter a lot for the conduct of MP)

    \item \textbf{The missing interaction between literatures} --- Economics and political science study CBI without citing each other. The economics community is institutionally embedded; the political science critique has no equivalent footprint. This non-interaction is itself evidence.
\end{itemize}

% -------------------------------------------------------
\subsection{Regulatory Capture and Epistemic Foreclosure (\textasciitilde500 words)}

\textit{Establish the framing concept. This paper is not about acquisition-capture; it is about foreclosure --- how a literature removes a question from democratic contest by never posing it.}

\begin{itemize}
    \item \textbf{Dal B\'{o} (2006): broad definition} --- Capture includes monetary policy explicitly. Pre-empts the objection that this only applies to utility regulation.

    \item \textbf{Stigler-Peltzman black box} --- Stigler: regulation is acquired by the industry and designed primarily for its benefit. Peltzman extends this into a model. Neither explains the internal mechanism --- Dal B\'{o} opens the box.

    \item \textbf{Pivot away from acquisition-capture} --- I do not claim finance engineered the literature. The mechanism is not lobbying or bribery.

    \item \textbf{Kwak (2013): cultural capture} --- Regulators influenced through group identification and socialisation, not bribes. No explicit coordination required. A precedent for non-acquisition mechanisms.

    \item \textbf{Epistemic foreclosure (the framing concept)} --- A literature can remove a question from democratic contest by never posing it. The goals of monetary policy are foreclosed not by being defeated in debate but by being folded into the definition of ``independence,'' so that contesting them registers as an attack on independence itself.

    \begin{itemize}
        \item Wansleben (2020) is important here
    \end{itemize}

    \item \textbf{GAO reports (2017, 2019)} --- Official acknowledgement of capture risk at the Fed and FDIC.

    \item \textbf{Posen (1995)} --- ``National differences in the degree of central bank independence in the postwar period resulted instead from differences in financial sector opposition to inflation'' (p.\ 253). The foreclosed default is the creditor-friendly one. \textit{Scope: noted, not causal.}
\end{itemize}

% -------------------------------------------------------
\section{Methodology (\textasciitilde950 words)}

\subsection{Research Design (\textasciitilde300 words)}

\begin{itemize}
    \item \textbf{Theory-testing single case study: the Federal Reserve} --- The Fed is the appropriate case because of its outsized importance in the development of CBI globally and the economics literature more broadly. The dual mandate is directly referenced in the index critique.

    \begin{itemize}
        \item No ECB basically because of length constraints
    \end{itemize}

    \item \textbf{Process tracing (Beach and Pedersen, 2019)} --- ``A research method for tracing causal mechanisms using detailed, within-case empirical analysis of how a causal mechanism operated in real-world cases'' (p.\ 1); ``we shift the analytical focus from causes and outcomes to the hypothesized causal mechanisms in between'' (p.\ 1). Inferences are made using ``the correspondence between the hypothetical empirical fingerprints that might have been left by the operation of a mechanism and the actual empirical observations we find in a case'' (p.\ 3).
    \begin{itemize}
        \item \textbf{Bennett and Checkel (2015)} --- Complementary guidelines for applying process tracing in political science, particularly on using within-case evidence to evaluate competing causal explanations.
    \end{itemize}

    \item \textbf{A sequential causal chain} --- Three sections are causally linked legs of one mechanism: the epistemological prior (4.1) folds the goal into the definition of independence; those assumptions travel into policy and institutional decisions through a few channels (4.2); that transmission forecloses legislative revision (4.3). Each step is downstream of the prior one.
\end{itemize}

\subsection{Causal Chain and Evidentiary Tests (\textasciitilde400 words)}

\begin{itemize}
    \item \textbf{Step 1 --- hoop test: the literature folds the goal into the definition of independence} --- Necessary condition and origin of the chain. If the literature does not treat the price-stability goal as constitutive of independence, the whole mechanism fails at its source. \textit{Falsifiable: does the literature ever treat the price-stability goal as a contestable choice, or always as a fixed premise?}

    \item \textbf{Step 2 --- hoop test: assumptions transmitted through institutional vessels} --- Necessary condition and middle leg. The assumptions from Step 1 must travel from academic literature into institutional practice through specific, traceable channels.

    \item \textbf{Step 3 --- smoking gun: legislative revision foreclosed using the same assumptions} --- The chain closes here. Goal-revision framed as an attack on independence --- using the same definitional fusion that originated in Step 1 --- is the circularity that confirms the mechanism is operating end-to-end.

    \item \textbf{Hoop test (Van Evera 1997)} --- Necessary but not sufficient: absent fingerprint kills the hypothesis; present fingerprint keeps it alive but doesn't confirm it.

    \item \textbf{Smoking gun test (Van Evera 1997)} --- Sufficient but not necessary: if found, strongly confirms the hypothesis; its absence doesn't eliminate it. Decisive because the evidence would be very hard to explain under any rival hypothesis.

    \item \textbf{Rival hypothesis limitation} --- \textit{[TO REVISIT]} Technocratic legitimacy (Posen) can explain Step 3 in isolation but can't explain why the measurement apparatus in Step 1 folds the goal into the definition, or why the discourse in Step 2 reproduces it.
\end{itemize}

\subsection{Rival Hypothesis (\textasciitilde250 words)}

\begin{itemize}
    \item \textbf{Posen (PIIE 2025): technocratic legitimacy} --- CBI works; legislators defer because it is rational to defer given the Fed's track record.

    \item \textbf{Insufficient} --- Explains the outcome but not the definitional fusion, not the discourse patterns, not the redirected reform energy.

    \item \textbf{Infrastructural power (Wansleben)} --- The Fed is a bureaucracy, and its innovations in monetary policy (governing techniques) have allowed it to grow and prosper. On this account, financial interests are not the driver of this development. \textit{Engage honestly as the strongest alternative.}

    \item \textbf{The rival hypothesis is itself a product of the epistemic environment under scrutiny} --- Posen's argument reproduces the assumptions of the pro-CBI literature. Its inability to see the foreclosure is exactly what the foreclosure thesis predicts.
\end{itemize}

% -------------------------------------------------------
\section{Analysis (\textasciitilde5,750 words)}

\subsection{The Epistemological Prior: Modeling Starting Conditions (\textasciitilde1,800 words)}
\label{sec:epistemology}


\textit{Step 1: the foreclosure is built in. There is an entire literature ingraining a single normative default at the very beginning. The motivating question is how two literatures --- economics and political science --- came to such radically different conclusions, and how the conclusions of one came to be encoded in institutional design.
}

\subsubsection{The Conservative Central Banker}

\begin{itemize}

    \item \textbf{Inflation has heterogeneous distributional effects} --- but the models assume one social objective function, erasing the conflict before analysis begins.

    \item \textbf{Price stability as unexamined social optimum} --- Kydland-Prescott assume low inflation is consistent with the public's preferences. Every subsequent paper inherits it.

    \item \textbf{Rogoff's conservative central banker} --- A central banker is not supposed to reflect the interests of the public but to be more inflation-averse than the public. One preference is locked in as the default.

    \item \textbf{Posen and the interests of finance} --- Finance is more inflation averse than the public.

    \begin{itemize}

    \item \textbf{Implication} --- The locked-in default happens to resemble the interests of finance. \textit{Scope: noted, not causal.}

    \item \textbf{The third best solution} --- Rogoff's preferred solution assumes market employment is being suppressed by distortions like unemployment insurance. That is a contestable normative claim, not a neutral starting point.

    \end{itemize}

    \item \textbf{Payoff} --- The model assumes away the distributional politics before analysis begins. The default is normatively loaded, not neutral.

    \end{itemize}

\subsubsection{Empirical evidence on CBI}

\textit{Describe general strategy of literature; Measuring de jure independence, building indices, etc... The foreclosure is demonstrated in the measurement apparatus itself.}

\begin{itemize}

    \item \textbf{De jure measurement conflates mandate content with operational autonomy} --- The indices measure what the mandate contains, not how independently it is pursued. These are categorically different things.

    \item \textbf{The circularity} --- Independence is defined partly by the goal. The objectives component scores price-stability mandates as \textit{more} independent, so a dual mandate scores as \textit{less} independent. Cukierman ``conservative bias'' quote is the smoking gun.

    \item \textbf{Mandate revision toward employment lowers most index scores} --- If Congress exercised its constitutional authority to revise the mandate, the Fed's independence score would fall, though instrument independence is untouched. That proves the circularity by thought experiment.

    \item \textbf{US/Turkey anomaly} --- The US scores below Turkey on de jure indices. Turkey is Posen's own example of a country that lost independence and got inflation. The measure cannot be tracking real independence; it is tracking goal-conformity.

    \item \textbf{The free lunch claim} --- Alesina and Summers say lower inflation, no output cost. The trade-off Rogoff acknowledged is gone. Echoed through additional sources.

    \item \textbf{Conclusion} --- To contest the goal democratically is, by construction, to be scored as less independent. Foreclosure is demonstrated in the measurement apparatus itself.

    \end{itemize}
    \subsubsection{Non-interaction of literatures}
    \begin{itemize}

    \item \textbf{Two non-interacting communities demonstrated textually} --- Most modern work on CBI ignores the entirety of the political science literature. Example: Romelli

    \item \textbf{The incommensurability of starting questions} --- The economics literature asks how to optimise monetary policy. The political science literature asks whose interests monetary policy serves. These are not competing answers to the same question --- they are different questions entirely.

  %  \item \textbf{Schneider, Strahan and Yang (2023)} --- Capture framework applied to Fed stress testing in a peer-reviewed journal. The argument is established in the literature even before applying it to monetary policy.

\end{itemize}


% -------------------------------------------------------
\subsection{Vessels of Transmission (\textasciitilde1,600 words)}

\textit{Step 2: the foreclosure moves from journals into the institution. The economics literature prior doesn't stay in journals. It travels into institutions through specific channels --- and those channels are what this section traces.}

\subsubsection{How the Prior Travels: Governance, Communication, and Expertise (\textasciitilde800 words)}

\begin{itemize}

    \item Christopher Waller 2025 Speech

    \item \textbf{Governors reproduce the default as settled fact} --- Fed officials cite the exact literature (Meyer, Mishkin, Kohn, Bernanke, Evans citing Rogoff / Cukierman / Alesina). The ``special interest'' they name as the threat is always democratic / government pressure to overstimulate, never a contestable-goals argument.

    \item \textbf{The tell} --- Meyer (2000) and Fischer (2015), from inside the institution, name the dual-mandate-penalty convention explicitly --- proving the foreclosure is recognized and reproduced, not accidental.

    \item \textbf{Downey: Fed never fully public} --- Reserve bank boards are two-thirds banker or banker-appointees. The primary dealer system concentrates monetary transmission in a small set of major financial institutions. Financial interests are designed into the institution.

    \item \textbf{Bernanke (\textit{The Courage to Act}, 2015)} --- Congress didn't understand what QE was while it was happening. The Fed's own chair confirms Dal B\'{o}'s asymmetric information mechanism. The technocratic claim isn't just a justification --- it actively disables oversight.

    \begin{itemize}
        \item The Fed as a research institution matters a lot- need to develop this idea further
    \end{itemize}

    \item \textbf{Schnabel (2026)} --- ``The pandemic episode was, hence, a testament to monetary dominance. It underscored the strength of an institutional framework anchored in the primacy of price stability.'' Reproduces the epistemological prior in real time. Her financial dominance framing mirrors Lindblom's market as prison argument from the other direction --- both describe financial sector structural power constraining policy, but Schnabel presents it as a condition to manage rather than a form of foreclosure.

    \begin{itemize}
        \item She works for the ECB- notable caveat. Doubtful a Fed official would say this, given the dual mandate
    \end{itemize}

    \item \textbf{PIIE 2025 as institutional exhibit} --- The monetary policy mandate is treated as beyond democratic contestation. Capture risk is acknowledged, but only for supervisory functions. The prior is being reproduced in real time.

    \item \textbf{Adolph: greater legal independence amplifies shadow principal influence} --- More independent central bank = more powerful agent = more valuable to financial sector principals.

    \item \textbf{Rockefeller/Wriston historical asymmetry} --- Financial sector critique of the Fed has always targeted specific constraints, not independence itself. After Volcker, the critique flips to attacking looseness (Druckenmiller). The asymmetry runs historically deep --- it is structural, not incidental.
\end{itemize}

\subsubsection{Pre-established influence: Revolving Door}

\textit{The prior reproduces itself through personnel. The career pipeline from finance to central bank to finance ensures the epistemic community that produces the literature also staffs the institution. Scope: this is transmission of a worldview, not proof of finance intent.}

\begin{itemize}

    \item \textbf{ABA revolving door: Olson; Duke} --- The same individuals move between the ABA and the Fed Board of Governors. The epistemic community reproduces itself through appointments.

    \item \textbf{Bowman: community bank $\rightarrow$ Vice Chair for Supervision $\rightarrow$ SVB reframing} --- Adolph's shadow principals mechanism made concrete. Career trajectory predicts policy outcomes. SVB framed as a management failure, not a structural regulatory failure. Supervision staff cut 30\%. Post-SVB reforms reversed.

    \item \textbf{Wirsching (2018)} --- Finance-sector backgrounds empirically predict deregulatory outcomes. Pattern predicted by Adolph, explained by Kwak.

    \item \textbf{Federal Advisory Council} --- The four largest banks have a formal governance channel that demonstrably influences the Federal Funds rate. This isn't revolving door dynamics --- it is structural institutional design.

    \item \textbf{Kwak: cultural capture through socialisation} --- Regulators identify with regulated firms through repeated interaction. No coordination required. The prior reproduces through professional culture.
\end{itemize}

% -------------------------------------------------------
\subsection{Legislative Blockage (\textasciitilde2,350 words)}

\textit{Step 3: foreclosure completes; revision becomes unthinkable. The defences of legislative inaction use the same assumptions exposed in Section 4.1. That is the circularity: the definition that folds the goal into independence is the same definition mobilised to argue revision is illegitimate.}

\begin{itemize}
    \item \textbf{Congress holds goal-setting authority but the goal is now equated with independence} --- Powell concedes Congress sets the goals. Yet because the goal is fused with independence, any revision is framed as an attack on independence, and therefore politically radioactive.

    \item \textbf{Volcker 1979: the original encoding event} --- Full regime shift toward price stability with no legal change. The dual mandate was operationally overwritten without democratic revision. The default was installed before it was ever contested.

    \item \textbf{Trump / Turkey contrast} --- Overt attacks on instrument independence get resisted and named. But goal-revision never even reaches contest, because it is pre-coded as illegitimate. The two are categorically different.

    \item \textbf{Lindblom (1982): market as prison} --- Any attempt to revise the mandate triggers financial market volatility. No lobbying required. The structural deterrent grows as the financial sector grows --- UK banking sector went from roughly 100\% to 450\% of GDP since the 1970s. Schnabel's financial dominance framing confirms this mechanism from within the central banking literature itself.

    \item \textbf{Dodd-Frank as failed critical juncture} --- Post-2008 was the strongest reform moment in decades. Reform got channelled into adjustments --- Class A director reform, stress testing --- that foreclosed fundamental revision. QE had already recapitalised the financial sector (Jacobs and King 2018 --- ``political multiplier effect'').

    \item \textbf{McCubbins and Schwartz (1984): fire alarm oversight} --- Congress only intervenes when interests trigger the alarm. The alarm never gets pulled on mandate questions because contesting the mandate is pre-coded as politicizing the Fed.

    \item \textbf{Bernanke/asymmetric information} --- Congress can't revise what it doesn't understand. The technocratic claim does active legislative work --- it doesn't just justify independence, it disables the mechanism that would produce revision.

    \item \textbf{PIIE 2025 as foreclosure exhibit} --- The expert community advising Congress on monetary policy shares the same assumptions exposed in Section 4.1. The defences of independence are the prior in institutional form.

    \item \textbf{The circularity stated explicitly} --- The definition that folds goal into independence is the same definition mobilised to defend against legislative revision. The rival hypothesis isn't an independent check --- it's produced by the same epistemic environment.

    \item \textbf{Close} --- The system requires no coordination. The goal is protected because contesting it is definitionally ``politicizing the Fed.'' That is the foreclosure completed.
\end{itemize}

% -------------------------------------------------------
\section{Conclusion (\textasciitilde650 words)}

\begin{itemize}
    \item \textbf{Epistemic foreclosure as unifying mechanism} --- The prior folds the goal into the definition of independence $\rightarrow$ travels through personnel and institutional discourse $\rightarrow$ makes democratic revision of the goal structurally unthinkable. The cycle is self-reinforcing.

    \item \textbf{The ``ingrained'' framing} --- The normative default is not external to the Fed's legitimating framework. It was encoded in it from the beginning.

    \item \textbf{Contribution 1} --- Epistemic foreclosure as a mechanism, distinct from acquisition-capture; reframes the CBI debate from ``how much independence'' to ``independence from whom, for what goal.''

    \item \textbf{Contribution 2} --- CBI measurement is circular. Measurement is based on legal goals, not instrument independence.

    \item \textbf{Technocratic claim as mechanism of circularity} --- Expertise justifies insulation $\rightarrow$ insulation means the goal goes unchallenged $\rightarrow$ unchallenged, the goal validates the expertise. The loop closes.

    \item \textbf{Distributional payoff} --- This is not neutral. The locked-in default has winners (creditors) and losers, so foreclosure is a political act disguised as a technical one.

    \item \textbf{Scope honesty} --- This is a claim about the literature and its transmission, not a proven conspiracy of finance.

    \item \textbf{Caveat} --- This is not an argument for Trump's ability to remove Powell. That's a different problem --- authoritarianism and central banking --- and a separate research agenda. This argument is about democratic design, particularly on the ``goal'' side of CBI.

    \item \textbf{Closing} --- The question isn't whether the Fed should be independent from political interference. It's whether the goals of monetary policy can be democratically contested at all, given a literature that defines independence as adherence to one goal.
\end{itemize}

\vspace{1cm}
\noindent\textit{Total: \textasciitilde10,000 words}

\end{doublespace}

\end{document}